% From Ideas to Action: a detailed write-up of the model in ideas_to_action.pdf
% !TEX program = xelatex
% Compile with XeLaTeX:  latexmk -xelatex writeup.tex
\documentclass[11pt]{article}

\usepackage{amsmath,amsthm}
\usepackage{fontspec}
\usepackage{unicode-math}
\setmainfont{STIX Two Text}
\setmathfont{STIX Two Math}
\usepackage[margin=1.1in]{geometry}
\usepackage{booktabs,xcolor,microtype}
\usepackage[hidelinks]{hyperref}

\definecolor{nap}{HTML}{9A6A1F}
\definecolor{ber}{HTML}{2F5D8A}
\definecolor{lan}{HTML}{3D7357}
\definecolor{muted}{HTML}{5A5A5A}
\newcommand{\Nap}{\textcolor{nap}{Napoleon}}
\newcommand{\Ber}{\textcolor{ber}{Berthier}}
\newcommand{\Lan}{\textcolor{lan}{Lannes}}

\newtheorem{definition}{Definition}
\newtheorem{assumption}{Assumption}
\newtheorem{lemma}{Lemma}
\newtheorem{proposition}{Proposition}
\newtheorem{corollary}{Corollary}
\newtheorem*{definition*}{Definition}
\newtheorem*{lemma*}{Lemma}
\newtheorem*{theorem*}{Theorem}
\newtheorem*{proposition*}{Proposition}
\theoremstyle{remark}
\newtheorem*{remark}{Remark}

\newenvironment{intuition}{\begin{quote}\small\color{muted}\textit{Intuition.}\ }{\end{quote}}

\newcommand{\E}{\mathbb{E}}
\DeclareMathOperator{\tr}{tr}
\DeclareMathOperator{\Var}{Var}
\newcommand{\Ycal}{\mathcal{Y}}
\newcommand{\Rcal}{\mathcal{R}}
\newcommand{\Acal}{\mathcal{A}}
\newcommand{\Scal}{\mathcal{S}}
\newcommand{\accept}{\text{accept}}

\linespread{1.08}
\setlength{\parskip}{0.35em}

\title{From Ideas to Action:\\[0.2em] \large Strategy, Solving, and Execution in a Model of AI and Growth}
\author{}
\date{October 2026}

\begin{document}
\maketitle

\begin{abstract}
Producing anything takes three kinds of work: deciding what to do, working out how to do it, and doing it. We model these as three agents, named after Napoleon, his chief of staff Berthier, and his marshal Lannes. \Nap{} screens an endless supply of ideas using noisy signals and must decide which few to pursue; \Ber{} turns a chosen plan into a detailed solution by approximating the true problem with a finite family of functions; \Lan{} carries the solution out, but only with the actions that are physically feasible. We define one improvement for each agent (better ideation, better problem solving, and better physical execution) and study all seven combinations, among them AI (better problem solving), RSI (adding better ideation), and AI with robotics (adding better execution). One principle organizes the results: \emph{the agent left behind sets the limit}. Any single improvement raises output to a ceiling set by another agent. Any pair raises it further, to a ceiling set by the third; with thin-tailed idea values, no pair sustains growth. Improving all three sustains growth, at the pace of the slowest improvement. The physical frontier also interacts with ideas: payoff is supermodular in a plan's ambition and the frontier, so better execution raises the return to better ideas. In a research sector, where \Ber's solution is itself the output and there is no \Lan, RSI is already enough.
\end{abstract}

\tableofcontents
\bigskip

%=====================================================================
\section{Introduction}
%=====================================================================

Suppose you want to know where the economic value of artificial intelligence will come from. One view is that AI will be extraordinarily good at \emph{solving} problems: give it a well-posed task and it will find a near-optimal answer quickly. Another view is that AI will eventually be good at \emph{deciding which problems are worth solving}. A third consideration is that, however good the thinking, something eventually has to be \emph{done} in the physical world, and that may be the binding constraint.

This note builds a small model that keeps these three kinds of work separate, and asks which of them can sustain growth. The model has three agents:
\begin{itemize}
  \item a \textbf{strategist}, \Nap, who decides \emph{what} to do;
  \item a \textbf{solver}, \Ber, who works out \emph{how};
  \item an \textbf{executor}, \Lan, who \emph{does} it.
\end{itemize}
The names come from the Napoleonic army: Napoleon decided what each campaign should achieve; Louis-Alexandre Berthier, his chief of staff, turned those intentions into detailed orders, routes, and timetables; Jean Lannes, one of his marshals, carried the orders out on the battlefield.

\begin{intuition}
The point of separating the three is that they fail in different ways. A strategist fails by pursuing the wrong idea. A solver fails by getting the details wrong for a good idea. An executor fails because the world does not allow exactly what the plan calls for. Improving one kind of work leaves the others' failures in place, and that is the source of every result below.
\end{intuition}

\paragraph{Results in brief.} We study the seven combinations of improvements in three groups.
\begin{enumerate}
  \item \textbf{One improvement has a ceiling} (Propositions~\ref{prop:ai}--\ref{prop:execution}). Better problem solving can at most remove its own errors and fill downstream capacity. Better ideation can at most sort plans perfectly. Better execution cannot plan faster. Each raises output once, to a ceiling set by another agent. Heavy-tailed idea values are the one way out, for faster screening (Section~\ref{sec:heavy}).
  \item \textbf{Two improvements still have a ceiling, set by the agent left behind} (Propositions~\ref{prop:frontier}--\ref{prop:ideationrobotics}). For RSI it is \Lan: the physical frontier bounds output whatever thinking does, and with thin tails RSI grows at most like $\sqrt t$. For AI with robotics it is \Nap's screening speed. For ideation with robotics it is \Ber's solving speed.
  \item \textbf{Ambition and the frontier are complements} (Lemma~\ref{lem:complements}). The gain from moving the physical frontier is largest for the most ambitious plans, and the best plans sit just beyond the frontier.
  \item \textbf{All three improvements sustain growth} (Proposition~\ref{prop:allthree}), at the rate of the slowest improvement, even with thin tails.
  \item \textbf{Research is different} (Propositions~\ref{prop:resceiling} and~\ref{prop:resgrowth}). With no executor, AI alone lifts research to \Nap's ceiling, and RSI sustains growth.
\end{enumerate}
Section~\ref{sec:experiment} describes a numerical experiment that runs all seven combinations from the same starting economy.

\paragraph{Related literature.} The model builds on idea-based growth (Romer 1990; Jones 1995; Weitzman 1998) and on research as search over a distribution of ideas (Kortum 1997), where, as in our Propositions~\ref{prop:thin} and~\ref{prop:ideationrobotics}, sustained growth from search needs heavy-tailed ideas. On AI and growth, Aghion, Jones and Jones (2019) emphasize bottlenecks and Acemoglu and Restrepo (2018) tasks. Separating problem solving from production echoes knowledge hierarchies (Garicano 2000), and \Nap's screening problem is one of sequential information acquisition (Wald 1947). Our contribution is to separate judgment, problem solving, and physical execution, and to ask which combination of improvements sustains growth.

%=====================================================================
\section{The environment}
%=====================================================================

\subsection{Plans, actions, and payoffs}

A \textbf{plan} $\xi\in\Xi$ says what to do: an objective and a broad course of action, without operational detail. A plan is carried out by a detailed \textbf{action} $a\in\Acal(\xi)\subseteq\mathbb{R}^d$; the action specifies the details (orders, routes, timing, allocation). Taking action $a$ in pursuit of plan $\xi$ yields \textbf{payoff} $W(a;\xi)$.

\begin{assumption}[Smooth problems]\label{ass:smooth}
For every plan, $\Acal(\xi)$ is compact and $W(\cdot\,;\xi)$ is infinitely differentiable ($C^\infty$) in $a$.
\end{assumption}

The \textbf{ideal solution} and the \textbf{value} of a plan are
\begin{equation}\label{eq:ideal}
a^*(\xi)\in\arg\max_{a\in\Acal(\xi)}W(a;\xi),\qquad v(\xi)=W\big(a^*(\xi);\xi\big).
\end{equation}
Compactness and continuity guarantee that the maximum exists. The value $v(\xi)$ is what the plan would deliver if every detail were worked out perfectly and every detail could be carried out exactly.

\begin{intuition}
The plan is the ``what''; the action is the ``how''. Two plans can share the same space of actions but rank them very differently. The value $v(\xi)$ is a best case that no one actually achieves: it is the yardstick against which the solver's and executor's shortfalls are measured.
\end{intuition}

\subsection{How a plan becomes output}

Every plan that is pursued passes through three stages:
\[
\underset{\text{\Nap's plan}}{\xi}\ \xrightarrow{\ \text{\Ber}\ }\ \underset{\text{solution}}{\hat a}\ \xrightarrow{\ \text{\Lan}\ }\ \underset{\text{action}}{\tilde a}\ \longrightarrow\ \underset{\text{output}}{Y=W(\tilde a;\xi)}.
\]
\Nap{} chooses $\xi$; \Ber{} produces a detailed solution $\hat a$; \Lan{} carries out an action $\tilde a$, which need not equal $\hat a$. Output from the project is the payoff of what actually gets done, $Y=W(\tilde a;\xi)$.

Each agent comes in many copies working in parallel: $n_A$ Napoleons, $n_B$ Berthiers, and $n_L$ Lanneses. Each stage takes time. We describe the downstream stages first, because \Nap's choice depends on what happens to a plan after he chooses it, much as one solves a dynamic program backwards.

%=====================================================================
\section{\texorpdfstring{The executor: \Lan}{The executor: Lannes}}
%=====================================================================

\setcounter{definition}{2}\begin{definition}[Executor, or \Lan]\label{def:lannes}
An executor, or a \Lan, is a pair $L=(\tau_L,\Rcal)$:
\begin{itemize}
  \item the \textbf{budget set} $\Rcal\subseteq\Acal(\xi)$ is the set of actions that can physically be carried out;
  \item the \textbf{execution time} $\tau_L$ is how long it takes to carry out one plan. With $n_L$ executors working in parallel, plans are carried out at rate $\mu_L=n_L/\tau_L$.
\end{itemize}
\end{definition}

\paragraph{Problem L (execution).} \Lan{} receives \Ber's solution $\hat a$ and carries out the feasible action closest to it:
\begin{equation}\label{eq:lannes}
\tilde a=\arg\min_{r\in\Rcal}\ \lVert r-\hat a\rVert.
\end{equation}

\paragraph{Execution loss.} What \Lan{} loses, relative to doing exactly what \Ber{} asked, is
\begin{equation}\label{eq:kappaL}
\kappa_L(\xi)=W(\hat a;\xi)-W(\tilde a;\xi).
\end{equation}
If $\hat a\in\Rcal$, then $\tilde a=\hat a$ and $\kappa_L=0$.

\paragraph{Best feasible value.} Because the executed action always lies in $\Rcal$, no plan can ever deliver more than
\begin{equation}\label{eq:Ubar}
\bar U(\xi)=\max_{r\in\Rcal}W(r;\xi)\ \le\ v(\xi).
\end{equation}
The inequality holds because $\Rcal\subseteq\Acal(\xi)$ and $v(\xi)$ is the maximum over $\Acal(\xi)$. The quantity $\bar U(\xi)$ depends only on the plan and on the budget set; it does not depend on \Ber.

\begin{intuition}
The executor's difficulty is not that he is clumsy. It is that the world offers a limited menu. A detailed solution can call for something that no available technology can do; the executor then does the closest thing he can. Notice that $\kappa_L$ can be negative: if \Ber's solution is off in a direction that the budget set happens to correct, the closest feasible action may be \emph{better} than what \Ber{} proposed. That is why $\kappa_L$ is not written as an absolute value: it is the signed shortfall, and the losses below have to add up.
\end{intuition}

%=====================================================================
\section{\texorpdfstring{The solver: \Ber}{The solver: Berthier}}
%=====================================================================

\Ber{} must turn a vague plan into a concrete action. The true problem $\max_aW(a;\xi)$ is usually too complicated to solve exactly, so \Ber{} replaces it with a tractable approximation and solves that.

\setcounter{definition}{1}\begin{definition}[Solver, or \Ber]\label{def:berthier}
A solver, or a \Ber, is a pair $B=(\tau_B,k)$. \Ber{} replaces the true payoff by its best approximation from a \textbf{sieve} $\mathcal{W}_k$, a family of functions of complexity $k$:
\begin{equation}\label{eq:sieve}
\hat W(\cdot\,;\xi)\in\arg\min_{f\in\mathcal{W}_k}\ \sup_{a\in\Acal(\xi)}\big|f(a)-W(a;\xi)\big|.
\end{equation}
\begin{itemize}
  \item The \textbf{sieve size} $k$ measures how rich a model \Ber{} can handle (for example, the degree of a polynomial, the number of spline knots, or the size of a neural network).
  \item The \textbf{solving time} $\tau_B$ is how long it takes to solve one plan. With $n_B$ solvers, plans are solved at rate $\mu_B=n_B/\tau_B$.
\end{itemize}
\end{definition}

Examples of sieves are polynomials of degree at most $p$ (so $k\asymp p^d$ coefficients), splines with $k$ basis functions, or neural networks with $k$ parameters. The approximation in~\eqref{eq:sieve} is \emph{global}: it is chosen to be close to $W$ over the whole action set, so it does not require knowing in advance where the optimum lies.

\paragraph{Problem B (approximate, then solve).} \Ber{} solves the approximate problem exactly:
\begin{equation}\label{eq:berthier}
\hat{\mathcal{P}}(\xi):\qquad \hat a\in\arg\max_{a\in\Acal(\xi)}\hat W(a;\xi).
\end{equation}

\paragraph{Solution loss.} The cost of working with the approximation rather than the truth is
\begin{equation}\label{eq:kappaB}
\kappa_B(\xi)=\big|v(\xi)-W(\hat a;\xi)\big|=v(\xi)-W(\hat a;\xi).
\end{equation}
The absolute value is automatic: $\hat a\in\Acal(\xi)$ and $v(\xi)$ is the maximum of $W$ over $\Acal(\xi)$, so $W(\hat a;\xi)\le v(\xi)$.

\subsection{How good is the approximation?}

Define the \textbf{approximation error}, the distance from $W$ to the sieve:
\begin{equation}\label{eq:epsB}
\varepsilon_B(k)=\min_{f\in\mathcal{W}_k}\ \sup_{a\in\Acal(\xi)}\big|f(a)-W(a;\xi)\big|.
\end{equation}
Approximation theory gives rates for standard sieves. For polynomial or spline sieves on a compact set in $\mathbb{R}^d$ and a payoff in $C^s$, Jackson-type theorems give
\begin{equation}\label{eq:jackson}
\varepsilon_B(k)\ \le\ C_s\,k^{-s/d},
\end{equation}
where $C_s$ depends on bounds on the first $s$ derivatives of $W$ and on the domain. Because $W$ is $C^\infty$ (Assumption~\ref{ass:smooth}), \eqref{eq:jackson} holds for every $s$.

\begin{intuition}
A richer sieve fits the true payoff better, and the rate depends on two things. Smoother problems (larger $s$) are easier: there is less wiggle to capture. Problems with more decision variables (larger $d$) are harder: the exponent $s/d$ shrinks, so it takes many more terms to reach the same accuracy. This is the curse of dimensionality, and it is where a solver with a much larger sieve, such as a large model, should matter most: detailed problems with many interacting choices.
\end{intuition}

\subsection{From a good approximation to a good solution}

\begin{lemma}[Approximation bound]\label{lem:approx}
If $\sup_a|\hat W(a;\xi)-W(a;\xi)|\le\varepsilon$, then $\kappa_B(\xi)\le2\varepsilon$. In particular $\kappa_B(\xi)\le2\varepsilon_B(k)\le2C_sk^{-s/d}$.
\end{lemma}

\begin{proof}
Three inequalities, each using one fact:
\begin{align*}
W(\hat a;\xi)
&\ \ge\ \hat W(\hat a;\xi)-\varepsilon &&\text{($\hat W$ is within $\varepsilon$ of $W$ at $\hat a$)}\\
&\ \ge\ \hat W\big(a^*(\xi);\xi\big)-\varepsilon &&\text{($\hat a$ maximizes $\hat W$)}\\
&\ \ge\ W\big(a^*(\xi);\xi\big)-2\varepsilon &&\text{($\hat W$ is within $\varepsilon$ of $W$ at $a^*$)}.
\end{align*}
So $W(\hat a;\xi)\ge v(\xi)-2\varepsilon$, that is $\kappa_B(\xi)\le2\varepsilon$. The second statement follows from the definition~\eqref{eq:epsB} and the rate~\eqref{eq:jackson}.
\end{proof}

\begin{intuition}
If your map of the terrain is never off by more than $\varepsilon$ anywhere, then the peak you find on the map is at most $2\varepsilon$ below the true peak: you can be fooled by $\varepsilon$ at your chosen point and by $\varepsilon$ at the true peak, but by no more. The factor of two is all that separates ``my model is accurate'' from ``my decision is good''. Note what \Ber{} does \emph{not} get wrong: the goal. He takes \Nap's plan as given; his errors are all in the details.
\end{intuition}

%=====================================================================
\section{What a plan is worth}
%=====================================================================

Combining the two downstream stages, the payoff of a plan is
\[
Y=W(\tilde a;\xi)=v(\xi)-\big[v(\xi)-W(\hat a;\xi)\big]-\big[W(\hat a;\xi)-W(\tilde a;\xi)\big]=v(\xi)-\kappa_B(\xi)-\kappa_L(\xi).
\]
This identity is exact: it simply adds and subtracts $W(\hat a;\xi)$. We call the result the plan's \textbf{expected value},
\begin{equation}\label{eq:U}
U=u(\xi)=\E\big[W(\tilde a;\xi)\mid\xi\big]=\textcolor{nap}{v(\xi)}-\textcolor{ber}{\kappa_B(\xi)}-\textcolor{lan}{\kappa_L(\xi)}.
\end{equation}
In the model as written, $\hat a$ and $\tilde a$ are deterministic given $\xi$, so the expectation is trivial and $U=Y$. We keep the notation $\E[\,\cdot\mid\xi]$ because nothing below changes if either stage adds noise.

\begin{intuition}
A plan's worth is its ideal value, minus what is lost in solving it, minus what is lost in doing it. This is the number the strategist should care about. A brilliant idea that cannot be carried out, or that is too complicated to work out in detail, is not worth much.
\end{intuition}

\paragraph{Where plans come from.} Candidate plans are drawn independently from a pool, $\xi\sim F$. Each draw comes with its \emph{whole} payoff function $W(\cdot\,;\xi)$, and from it we compute $U=u(\xi)$. The distribution of $U$, which we call $G$, is therefore \emph{induced}: it depends on the pool $F$, and on how well \Ber{} and \Lan{} handle each plan. In particular, $G$ shifts whenever \Ber{} or \Lan{} improves. In general $G$ can be any distribution; Section~\ref{sec:gauss} gives conditions under which it is normal.

\subsection{A local view of the losses}

It helps to see the losses near the ideal solution. Expanding $W$ to second order around $a^*(\xi)$, where the gradient vanishes (if $a^*$ is interior), with $d=a-a^*(\xi)$,
\[
W(a;\xi)\approx v(\xi)-\tfrac12\,d^{\top}Hd,\qquad H=-\nabla_a^2W\big(a^*(\xi);\xi\big).
\]
If \Ber's approximation places the optimum at $\hat a=a^*(\xi)+\delta_B$, his loss is approximately $\kappa_B\approx\tfrac12\delta_B^{\top}H\delta_B$: quadratic in how far off he is, weighted by how sharply the payoff falls away from its peak. \Lan's loss is then the change in payoff from moving $\hat a$ to the nearest point of $\Rcal$.

\begin{intuition}
Solver errors come from approximation; executor errors come from infeasibility. Near the optimum, small errors are cheap (the loss is quadratic), which is why a moderately good solver already captures most of the value of a plan, and why improving him further has diminishing returns.
\end{intuition}

%=====================================================================
\section{\texorpdfstring{The strategist: \Nap}{The strategist: Napoleon}}
%=====================================================================

\subsection{Signals and policies}

\Nap{} faces an effectively unlimited supply of candidate plans, but he does not know any plan's value. What is scarce is the time to find out.

\setcounter{definition}{0}\begin{definition}[Strategist, or \Nap]\label{def:napoleon}
A strategist, or a \Nap, is a pair $N=(\tau_s,\Scal)$:
\begin{itemize}
  \item the \textbf{signal technology} $\Scal=\{P(\cdot\mid u)\}_u$ is a family of distributions on a signal space $\Sigma$. Given a plan with expected value $U=u$, signals $s_1,s_2,\dots\in\Sigma$ are drawn i.i.d.\ from $P(\cdot\mid u)$;
  \item the \textbf{signal time} $\tau_s$ is how long each signal takes.
\end{itemize}
\end{definition}

The process for one \Nap{} is a loop:
\begin{enumerate}
  \item draw a fresh candidate plan $\xi\sim F$, whose value $U$ is unknown;
  \item draw signals about it, one at a time, each taking time $\tau_s$;
  \item after each signal, decide to \emph{accept} the plan (send it to \Ber), \emph{reject} it (discard it for good), or \emph{wait} for another signal;
  \item after accepting or rejecting, return to step~1.
\end{enumerate}
Write $s_{1:n}=(s_1,\dots,s_n)$.

\begin{definition*}[Policy]
A \textbf{policy} is a sequence of decision rules
\[
\pi=(\pi_n)_{n\ge1},\qquad\pi_n:\Sigma^n\to\{\accept,\ \text{reject},\ \text{wait}\}.
\]
It stops after $N_\pi=\min\{n:\pi_n(s_{1:n})\ne\text{wait}\}$ signals. Let $\Pi$ be the set of policies with $\E[N_\pi]<\infty$.
\end{definition*}

\begin{intuition}
Think of a venture capitalist facing an endless stream of pitches, with limited time for due diligence on each and a fund that can back only a few companies a year. Each signal is a meeting, a reference check, a test. The policy says when he has learned enough to say yes or no. Waiting buys information about the current idea, but leaves less time for the next one.
\end{intuition}

\subsection{\texorpdfstring{How many plans are accepted: computing $Q(\pi)$}{How many plans are accepted: computing Q}}

Three quantities determine how many plans flow out of \Nap's office per unit of time.

\paragraph{The acceptance probability.} Let $P_u^n(ds_{1:n})=\prod_{k=1}^nP(ds_k\mid u)$ be the distribution of the first $n$ signals for a plan of value $u$. \Nap{} accepts at exactly the $n$-th signal on the set of histories
\[
A_n=\big\{s_{1:n}:\ \pi_j(s_{1:j})=\text{wait for all }j<n,\ \ \pi_n(s_{1:n})=\accept\big\}.
\]
The events ``accept at signal $n$'' are disjoint across $n$, so the chance that a plan of value $u$ is accepted, and the unconditional chance, are
\begin{equation}\label{eq:paccept}
P_\pi(\accept\mid u)=\sum_{n\ge1}P_u^n(A_n),\qquad P_\pi(\accept)=\int P_\pi(\accept\mid u)\,dG(u).
\end{equation}
A sharper signal technology makes $P_\pi(\accept\mid u)$ rise more steeply in $u$: good plans get accepted and bad ones rejected more reliably.

\paragraph{The length of an investigation.} \Nap{} is still waiting after $m$ signals on
\[
W_m=\big\{s_{1:m}:\ \pi_j(s_{1:j})=\text{wait for all }j\le m\big\},\qquad W_0=\text{all histories}.
\]
He draws an $n$-th signal exactly when his first $n-1$ signals lie in $W_{n-1}$. Using $\E[N]=\sum_{n\ge1}P(N\ge n)$ for a positive integer random variable,
\begin{equation}\label{eq:EN}
\E_\pi[N_\pi]=\sum_{n\ge1}P_\pi(N_\pi\ge n)=\sum_{n\ge1}\int P_u^{\,n-1}(W_{n-1})\,dG(u).
\end{equation}
One investigation takes $\tau_sN_\pi$ units of time, and $\tau_s\E_\pi[N_\pi]$ on average.

\paragraph{The renewal--reward theorem.} We need one result from probability to turn these per-plan quantities into rates.

\begin{theorem*}[Renewal--reward]
Let a process split into cycles $i=1,2,\dots$, where cycle $i$ lasts $T_i>0$ and earns reward $R_i$, with $(T_i,R_i)$ i.i.d.\ and $\E[T]<\infty$, $\E|R|<\infty$. Then, almost surely,
\[
\lim_{t\to\infty}\frac{\text{total reward earned by time }t}{t}=\frac{\E[R]}{\E[T]}.
\]
\end{theorem*}

\begin{proof}
Write $S_n=\sum_{i\le n}T_i$ and let $M(t)=\max\{n\ge0:\ S_n\le t\}$ be the number of cycles completed by time $t$. Suppose first that each cycle's reward is credited when the cycle ends.
\begin{enumerate}
  \item \emph{Cycles keep finishing.} By the strong law of large numbers, $S_n/n\to\E[T]\in(0,\infty)$ almost surely. So $S_n\to\infty$, which gives $M(t)<\infty$ for every $t$, and $S_n<\infty$ for every $n$, which gives $M(t)\to\infty$ as $t\to\infty$.
  \item \emph{The rate of completions.} By definition $S_{M(t)}\le t<S_{M(t)+1}$. Dividing by $M(t)$, which is positive for $t$ large,
  \[
  \frac{S_{M(t)}}{M(t)}\ \le\ \frac{t}{M(t)}\ <\ \frac{S_{M(t)+1}}{M(t)+1}\cdot\frac{M(t)+1}{M(t)}.
  \]
  Both outer terms tend to $\E[T]$ because $M(t)\to\infty$, so $M(t)/t\to1/\E[T]$.
  \item \emph{The rate of reward.} The reward credited by time $t$ is $C(t)=\sum_{i\le M(t)}R_i$, and
  \[
  \frac{C(t)}{t}=\Big(\frac{1}{M(t)}\sum_{i\le M(t)}R_i\Big)\cdot\frac{M(t)}{t}\ \longrightarrow\ \E[R]\cdot\frac{1}{\E[T]},
  \]
  by the strong law applied to the $R_i$ (which needs $\E|R|<\infty$) and step 2.
  \item \emph{Rewards earned during a cycle.} If reward accrues during a cycle, and the amount earned so far in cycle $i$ is at most $D_i$ in absolute value, with $(T_i,R_i,D_i)$ i.i.d.\ and $\E[D]<\infty$, then the reward earned by time $t$ differs from $C(t)$ by at most $D_{M(t)+1}$. The strong law gives $\frac1n\sum_{i\le n}D_i\to\E[D]$, hence $D_n/n\to0$, so
  \[
  \frac{D_{M(t)+1}}{t}=\frac{D_{M(t)+1}}{M(t)+1}\cdot\frac{M(t)+1}{t}\ \longrightarrow\ 0\cdot\frac{1}{\E[T]}=0. \qedhere
  \]
\end{enumerate}
\end{proof}

In our application every reward (a plan examined, a plan accepted, and the value of an accepted plan) is realized when the investigation ends, so step 4 is not needed.

\begin{intuition}
Over a long time $t$, about $t/\E[T]$ cycles fit, and each earns $\E[R]$ on average. Note that the rate is a ratio of averages, $\E[R]/\E[T]$, not an average of ratios. If half of all investigations take one signal and half take five, \Nap{} examines plans at rate $1/(3\tau_s)$, not at the average of $1/\tau_s$ and $1/(5\tau_s)$.
\end{intuition}

\paragraph{Putting it together.} Apply the renewal--reward theorem to one \Nap, with one cycle being one investigation, so $T_i=\tau_sN_i$:
\begin{itemize}
  \item with reward $1$ per cycle, plans are examined at rate $1/(\tau_s\E_\pi[N_\pi])$;
  \item with reward $\mathbf 1\{\accept\}$, plans are accepted at rate $P_\pi(\accept)/(\tau_s\E_\pi[N_\pi])$;
  \item with reward $U_i\mathbf 1\{\accept\}$, value is produced at rate $\E[U\mathbf 1\{\accept\}]/(\tau_s\E_\pi[N_\pi])$.
\end{itemize}
With $n_A$ strategists working in parallel, the \textbf{accepted flow} is
\begin{equation}\label{eq:Q}
Q(\pi)=\underset{\text{strategists}}{n_A}\times\underset{\text{plans examined by each}}{\frac{1}{\tau_s\,\E_\pi[N_\pi]}}\times\underset{\text{share accepted}}{P_\pi(\accept)}.
\end{equation}

\begin{intuition}
$Q(\pi)$ is a \emph{rate}, not a probability. Read it as: the strategists can process $n_A/\tau_s$ signals per unit of time; each plan uses $\E_\pi[N_\pi]$ signals on average, so they examine $(n_A/\tau_s)/\E_\pi[N_\pi]$ plans per unit of time; a share $P_\pi(\accept)$ of those is accepted.
\end{intuition}

\paragraph{The value of accepted plans.} By Bayes' rule,
\begin{equation}\label{eq:Uaccept}
\E_\pi[U\mid\accept]=\frac{\int u\,P_\pi(\accept\mid u)\,dG(u)}{P_\pi(\accept)},
\end{equation}
and by iterated expectations $\E_\pi[Y\mid\accept]=\E_\pi[U\mid\accept]$. Output per unit of time is the accepted flow times the average value of an accepted plan:
\begin{equation}\label{eq:output}
\Ycal(\pi)=Q(\pi)\,\E_\pi[U\mid\accept]=\frac{n_A}{\tau_s\,\E_\pi[N_\pi]}\int u\,P_\pi(\accept\mid u)\,dG(u).
\end{equation}

\subsection{\texorpdfstring{Capacity and \Nap's problem}{Capacity and Napoleon's problem}}

Accepted plans must be solved by \Ber{} and carried out by \Lan. Downstream, plans can be processed at most at the rate of the slower stage:
\begin{equation}\label{eq:K}
K=\min\{\mu_B,\ \mu_L\}=\min\Big\{\frac{n_B}{\tau_B},\ \frac{n_L}{\tau_L}\Big\}.
\end{equation}

\paragraph{Problem N (investigation).}
\begin{equation}\label{eq:problemN}
\max_{\pi\in\Pi}\ Q(\pi)\,\E_\pi[U\mid\accept]\qquad\text{subject to}\qquad Q(\pi)\le K.
\end{equation}
Dividing the constraint by the examination rate gives an equivalent per-plan form:
\[
P_\pi(\accept)\ \le\ \frac{K\,\tau_s\,\E_\pi[N_\pi]}{n_A}.
\]
\Nap{} may accept at most the share of plans that downstream capacity can absorb.

\begin{intuition}
Why would \Nap{} ever reject a plan? For two reasons. First, capacity is scarce: every plan he accepts takes a slot that a better plan could have used, so a merely decent plan is rejected because something better will come along. Second, some plans lose value: $U$ is net of the solver's and executor's losses, so a typical idea can be worth less than nothing once it is actually worked out and carried out. If capacity were unlimited and every plan had positive value, he would accept everything, and he would have no reason to investigate at all. Screening is valuable precisely because capacity is scarce and many ideas are bad.
\end{intuition}

%=====================================================================
\section{A worked example: Gaussian plans and signals}\label{sec:gauss}
%=====================================================================

This section derives closed forms for everything above under conditions that make $U$ normal.

\subsection{\texorpdfstring{When is $U$ normal?}{When is U normal?}}

$U$ is a derived quantity, so its distribution must come from the primitives. Suppose plans differ only in their \emph{level}:
\begin{equation}\label{eq:level}
W(a;\xi)=v(\xi)+w(a),\qquad v(\xi)\sim N(m,\sigma^2),
\end{equation}
with a common action set $\Acal$ and a common shape $w$ normalized so that $\max_aw(a)=0$.

\begin{lemma*}[Gaussian plans]
Under~\eqref{eq:level}, if the sieve $\mathcal{W}_k$ is closed under adding constants, then $\hat a$ and $\tilde a$ are the same for every plan, both losses are constants, and
\[
U=v(\xi)-\kappa_B-\kappa_L\ \sim\ N\big(m-\kappa,\ \sigma^2\big),\qquad\kappa=\kappa_B+\kappa_L.
\]
\end{lemma*}

\begin{proof}
If $f$ is a best approximation to $w$ from the sieve, then $v(\xi)+f$ is a best approximation to $v(\xi)+w$, because $\sup_a|(v+f)-(v+w)|=\sup_a|f-w|$ and adding a constant maps the sieve into itself. So $\hat W(\cdot;\xi)=v(\xi)+f$, whose maximizer $\hat a$ does not depend on $v(\xi)$. Problem L depends only on $\hat a$ and $\Rcal$, so $\tilde a$ does not depend on $\xi$ either. Then $\kappa_B=v(\xi)-W(\hat a;\xi)=-w(\hat a)$ and $\kappa_L=w(\hat a)-w(\tilde a)$ are constants, and $U=v(\xi)+w(\tilde a)=v(\xi)-\kappa_B-\kappa_L$ is a shift of a normal variable.
\end{proof}

\begin{intuition}
If all plans look alike in their details and differ only in how much they are worth, then the solver and executor handle them all identically, so they lose the same amount on each. The only thing left to screen on is the level, and that is normal by assumption. Outside these conditions $U$ need not be normal; for example, in Section~\ref{sec:physical}, more ambitious plans lose more in execution, which thins the upper tail of $U$.
\end{intuition}

Write $m_U=m-\kappa$ and $\sigma_U=\sigma$ from now on.

\subsection{Signals and the posterior mean}

Let signals be normal around the plan's expected value:
\[
s_k=U+\eta_k,\qquad\eta_k\sim N(0,\sigma_\eta^2)\ \text{i.i.d.},\qquad r=\sigma_\eta^2/\sigma_U^2.
\]
The ratio $r$ is the noise-to-signal ratio of a single signal.

\paragraph{Derivation of the posterior.} With a normal prior and normal signals, the posterior of $U$ after $n$ signals is normal. Its precision (inverse variance) is the sum of the prior precision and the signals' precisions, and its mean is the precision-weighted average of the prior mean and the signal average $\bar s_n=\frac1n\sum_{k\le n}s_k$:
\[
\mu_n=\E[U\mid s_{1:n}]=\frac{\sigma_U^{-2}m_U+n\sigma_\eta^{-2}\bar s_n}{\sigma_U^{-2}+n\sigma_\eta^{-2}}=m_U+\rho_n(\bar s_n-m_U),\qquad\rho_n=\frac{n\sigma_U^2}{n\sigma_U^2+\sigma_\eta^2}=\frac{n}{n+r}.
\]

\paragraph{Distribution of the estimate.} Before signals arrive, $\bar s_n=U+\bar\eta_n\sim N(m_U,\ \sigma_U^2+\sigma_\eta^2/n)$. Hence $\mu_n$ is normal with mean $m_U$ and variance
\[
\Var(\mu_n)=\rho_n^2\Big(\sigma_U^2+\frac{\sigma_\eta^2}{n}\Big)=\rho_n^2\cdot\frac{\sigma_U^2}{\rho_n}=\rho_n\sigma_U^2.
\]
(The middle step uses $\sigma_U^2+\sigma_\eta^2/n=\sigma_U^2/\rho_n$.)

\begin{intuition}
$\rho_n\in(0,1)$ is the \emph{accuracy} of \Nap's estimate: the share of the true variation in plan values that his estimate captures. With one very noisy signal, $\rho_n$ is near zero and his estimates barely differ across plans; with many signals, $\rho_n$ approaches one and he can see the plans' true differences. Equivalently, $\sqrt{\rho_n}$ is the correlation between his estimate and the truth.
\end{intuition}

\subsection{Napoleon's speed--accuracy trade-off}

Restrict \Nap{} to policies that take exactly $n$ signals per plan and then accept if the estimate clears a cutoff, $\mu_n\ge z$. (General sequential policies can do better, but the bounds in Section~\ref{sec:results} hold for all policies.)

\paragraph{Step 1: how selective he can be.} Each plan takes $n\tau_s$, so the strategists examine
\[
R_n=\frac{n_A}{n\,\tau_s}
\]
plans per unit of time. Capacity lets $K$ through, so when capacity binds \Nap{} accepts the top share
\[
q_n=\frac{K}{R_n}=\frac{K\,n\,\tau_s}{n_A}
\]
of examined plans. He does so with a cutoff on his estimate, and never accepts a plan he expects to be worth less than nothing:
\begin{equation}\label{eq:cutoff}
\accept\iff\mu_n\ge z_n,\qquad z_n=\max\Big\{0,\ m_U+\sigma_U\sqrt{\rho_n}\,\Phi^{-1}(1-q_n)\Big\}.
\end{equation}
The second term in the max is the cutoff that makes exactly a share $q_n$ of estimates clear it, since $\mu_n\sim N(m_U,\rho_n\sigma_U^2)$. The zero is there because accepting a plan with negative expected value lowers output even if a slot is free.

\paragraph{Step 2: how good his picks are.} Because the acceptance decision depends only on the signals, and $\E[U\mid s_{1:n}]=\mu_n$, the law of iterated expectations gives
\[
\E[U\mid\accept]=\E\big[\mu_n\mid\mu_n\ge z_n\big].
\]
Write $\mu_n=m_U+\sigma_U\sqrt{\rho_n}\,X$ with $X\sim N(0,1)$. When capacity binds, accepting means $X\ge c$ with $c=\Phi^{-1}(1-q_n)$. For a standard normal, the truncated mean is the inverse Mills ratio:
\[
\E[X\mid X\ge c]=\frac{1}{1-\Phi(c)}\int_c^\infty x\,\phi(x)\,dx=\frac{\phi(c)}{1-\Phi(c)},
\]
using $\int_c^\infty x\phi(x)dx=\big[-\phi(x)\big]_c^\infty=\phi(c)$. With $1-\Phi(c)=q_n$ this is
\[
S(q_n)=\frac{\phi\big(\Phi^{-1}(1-q_n)\big)}{q_n},
\]
the mean of the top share $q_n$ of a standard normal, in standard deviations. For example $S(0.05)=2.06$, $S(0.15)=1.55$, $S(0.5)=0.80$, and $S(1)=0$. Therefore
\[
\E[U\mid\accept]=m_U+\sigma_U\sqrt{\rho_n}\;S(q_n).
\]

\paragraph{Step 3: output.} Output is the accepted flow, $K$, times the value of each accepted plan:
\begin{equation}\label{eq:tradeoff}
\Ycal(n)=K\Big(m_U+\sigma_U\,\sqrt{\rho_n}\;S(q_n)\Big).
\end{equation}
Each term has a meaning:
\begin{itemize}
  \item $K$: plans carried out per unit of time;
  \item $m_U$: the value of an average plan;
  \item $\sigma_U$: how much plans differ, which is what there is to gain from choosing well;
  \item $\sqrt{\rho_n}$: how much of that difference \Nap{} can see after $n$ signals;
  \item $S(q_n)$: how far above average, in standard deviations, the plans he picks are.
\end{itemize}
Formula~\eqref{eq:tradeoff} assumes capacity binds and the capacity cutoff is positive. Otherwise the cutoff is zero, fewer than $K$ plans are accepted, and output is $R_n\,\E[\mu_n\mathbf 1\{\mu_n\ge0\}]=R_n\big[m_U(1-\Phi(c_0))+\sigma_U\sqrt{\rho_n}\,\phi(c_0)\big]$ with $c_0=-m_U/(\sigma_U\sqrt{\rho_n})$.

\begin{intuition}
More signals per plan raise accuracy $\sqrt{\rho_n}$, but they also mean fewer plans examined, so \Nap{} must accept a larger share $q_n$ and the selectivity premium $S(q_n)$ falls. Too few signals and he cannot tell plans apart; too many and he sees too few plans to be choosy. The optimal $n$ balances the two.
\end{intuition}

\paragraph{Choosing $n$: a numerical example.} Take $m_U=-2.5$, $\sigma_U=4$, $r=4$, $\tau_s=0.05$, $n_A=1$, and $K=1$. Then:

\begin{center}
\renewcommand{\arraystretch}{1.15}
\begin{tabular}{@{}c r r r r r@{}}
\toprule
$n$ & $R_n$ & $q_n$ & $\sqrt{\rho_n}$ & $S(q_n)$ & $\Ycal(n)$\\
\midrule
1 & 20.0 & 0.05 & 0.45 & 2.06 & 1.19\\
2 & 10.0 & 0.10 & 0.58 & 1.76 & 1.55\\
\textbf{3} & \textbf{6.7} & \textbf{0.15} & \textbf{0.65} & \textbf{1.55} & \textbf{1.57}\\
4 & 5.0 & 0.20 & 0.71 & 1.40 & 1.46$^*$\\
5 & 4.0 & 0.25 & 0.75 & 1.27 & 1.34$^*$\\
\bottomrule
\end{tabular}
\end{center}
\noindent{\small $^*$The capacity cutoff is negative, so the cutoff is zero and fewer than $K$ plans are accepted; output uses the formula below~\eqref{eq:tradeoff}.}

With $n=3$, each plan takes $0.15$ years, so \Nap{} examines about $6.7$ plans a year and accepts the top $15\%$, with cutoff $z\approx-2.5+4(0.65)(1.04)\approx0.21$. That gives exactly one accepted plan a year, $Q=0.15/(3\times0.05)=1=K$, and output $1.57$.

%=====================================================================
\section{Improvements and scenarios}\label{sec:improvements}
%=====================================================================

Each improvement relaxes exactly one agent's constraint. A \textbf{path} of an improvement is a sequence of technologies over time $t=0,1,2,\dots$, each at least as good as the last in the sense of the definition.

\setcounter{definition}{3}\begin{definition}[Better ideation]\label{def:ideation}
\textbf{Better ideation} is an improvement to \Nap: it replaces $(\tau_s,\Scal)$ with $(\tau_s',\Scal')$ where $\tau_s'\le\tau_s$, or $\Scal'$ is Blackwell more informative than $\Scal$ (so $\Scal$ is a garbling of $\Scal'$).
\end{definition}

\begin{definition}[Better problem solving]\label{def:solving}
\textbf{Better problem solving} is an improvement to \Ber: it replaces $(\tau_B,k)$ with $(\tau_B',k')$ where $\tau_B'\le\tau_B$ or $k'\ge k$.
\end{definition}

\begin{definition}[Better physical execution]\label{def:execution}
\textbf{Better physical execution} is an improvement to \Lan: it replaces $(\tau_L,\Rcal)$ with $(\tau_L',\Rcal')$ where $\tau_L'\le\tau_L$ or $\Rcal'\supseteq\Rcal$.
\end{definition}

\begin{definition}[Scenarios]\label{def:scenarios}
A \textbf{scenario} is a non-empty set of improvements, applied together, with every other agent's technology held fixed. There are seven:
\begin{center}
\renewcommand{\arraystretch}{1.1}
\begin{tabular}{@{}l c c c@{}}
\toprule
& \textcolor{nap}{Ideation} & \textcolor{ber}{Problem solving} & \textcolor{lan}{Execution}\\
\midrule
\multicolumn{4}{@{}l}{\textit{One improvement}}\\
\quad Better ideation alone & $\checkmark$ & & \\
\quad \textbf{AI} (better problem solving alone) & & $\checkmark$ & \\
\quad Better execution alone & & & $\checkmark$\\
\multicolumn{4}{@{}l}{\textit{Two improvements}}\\
\quad \textbf{RSI} (recursive self-improvement) & $\checkmark$ & $\checkmark$ & \\
\quad \textbf{AI with robotics} & & $\checkmark$ & $\checkmark$\\
\quad Ideation with robotics & $\checkmark$ & & $\checkmark$\\
\multicolumn{4}{@{}l}{\textit{All three}}\\
\quad \textbf{RSI with robotics} & $\checkmark$ & $\checkmark$ & $\checkmark$\\
\bottomrule
\end{tabular}
\end{center}
\end{definition}

\begin{intuition}
Better problem solving means faster solving ($\mu_B$ rises, which matters only while \Ber{} is the bottleneck) or richer models ($\varepsilon_B$ falls, which tightens the bound on $\kappa_B$). Better ideation means screening plans faster or judging them more sharply. Better execution means more capacity to act, or a menu of feasible actions that reaches further. The seven scenarios differ in which agent, if any, is left behind, and that turns out to decide which ceiling applies.
\end{intuition}

%=====================================================================
\section{The numerical experiment}\label{sec:experiment}
%=====================================================================

Each group of results below is followed by a numerical illustration. This section says exactly how it is computed.

\paragraph{The economy at year 0.} Gaussian plans (Section~\ref{sec:gauss}) with $m=0$ and $\sigma=4$, and one strategist ($n_A=1$). By the Gaussian-plans lemma both losses are the same for every plan, so $U\sim N(m-\kappa_B-\kappa_L,\ \sigma^2)$. The other year-0 parameters, and how each improvement changes them, are:
\begin{center}
\renewcommand{\arraystretch}{1.12}
\begin{tabular}{@{}l l l r r@{}}
\toprule
Improvement & Parameter & Meaning & Year 0 & Per year\\
\midrule
\textcolor{nap}{Ideation} (\Nap) & $\tau_s$ & years per signal & $0.05$ & $\div\,1.1$\\
& $r$ & signal noise, $\sigma_\eta^2/\sigma^2$ & $4$ & $\div\,1.1$\\
\textcolor{ber}{Problem solving} (\Ber) & $\mu_B$ & plans solved per year & $1$ & $\times\,1.25$\\
& $\kappa_B$ & loss from solving, per plan & $1.5$ & $\div\,1.25$\\
\textcolor{lan}{Execution} (\Lan) & $\mu_L$ & plans carried out per year & $2$ & $\times\,1.25$\\
& $\kappa_L$ & loss from execution, per plan & $1$ & $\div\,1.25$\\
\bottomrule
\end{tabular}
\end{center}
The parameters $m$, $\sigma$, and $n_A$ never change. A falling $\kappa_B$ stands for a growing sieve (Lemma~\ref{lem:approx}); a falling $\kappa_L$ stands for a budget set that reaches closer to \Ber's solution.

\paragraph{The procedure.}
\begin{enumerate}
  \item \emph{Same start.} Every scenario starts from the year-0 economy, where $K=\min\{1,2\}=1$, $m_U=-2.5$, \Nap{} takes $n=3$ signals, and output is $1.57$ a year (Section~\ref{sec:gauss}).
  \item \emph{Improve only what the scenario names.} In year $t$ the improved agents' parameters equal their year-0 values times the yearly factor to the power $t$; all other parameters stay at year 0.
  \item \emph{\Nap{} re-optimizes every year.} Given that year's $K_t=\min\{\mu_{B,t},\mu_{L,t}\}$, $m_{U,t}=m-\kappa_{B,t}-\kappa_{L,t}$, $r_t$, and $\tau_{s,t}$, he chooses the number of signals $n\in\{1,2,\dots\}$ and the cutoff~\eqref{eq:cutoff} to maximize output, using~\eqref{eq:tradeoff} or the formula below it when the cutoff is zero.
  \item \emph{Record.} Output in years 0, 10, 20, and 60; the growth rate as the average annual rate over years 50--60, $(\Ycal_{60}/\Ycal_{50})^{1/10}-1$; and the bottleneck in year 60.
\end{enumerate}
The \textbf{bottleneck} is \Nap{} if he examines no more plans than capacity allows, $n_A/(n\tau_{s,t})\le K_t$, so that his screening limits the flow; otherwise it is whichever of \Ber{} and \Lan{} has the lower rate. All numbers are exact expected values: there is no simulation noise. The parameters are illustrative, not a forecast.

%=====================================================================
\section{Results}\label{sec:results}
%=====================================================================

\subsection{Assumptions}

\begin{assumption}\label{ass:results}
\begin{enumerate}
  \item[(a)] Each result fixes the agents whose improvement it leaves out: their technologies do not change over time.
  \item[(b)] $\E|v(\xi)|<\infty$ under $F$, and output starts positive.
  \item[(c)] Where stated, plan values have thin tails: the variable named in the result, either the best feasible value $\bar U(\xi)=\max_{r\in\Rcal}W(r;\xi)$ or the ideal value $v(\xi)$, is $\sigma$-sub-Gaussian around its mean. For $\bar U$ with mean $\bar m=\E[\bar U]$, this means
  \[
  \E\Big[e^{\theta(\bar U-\bar m)}\Big]\le e^{\theta^2\sigma^2/2}\qquad\text{for all }\theta\in\mathbb{R}.
  \]
  \item[(d)] Screening outpaces capacity: $n_A/\tau_s\ge\sqrt e\,\bar K$, where $\bar K$ is the capacity of the fixed downstream agent named in the result ($\mu_L$ in Propositions~\ref{prop:judgment} and~\ref{prop:thin}, $\mu_B$ in Propositions~\ref{prop:execution} and~\ref{prop:ideationrobotics}).
\end{enumerate}
The \textbf{output growth rate} is $g_{\Ycal}=\limsup_{t\to\infty}\frac1t\ln\Ycal_t$. Write $\bar R=n_A/\tau_s$ for the fastest possible rate of examining plans (every plan needs at least one signal).
\end{assumption}

\begin{intuition}
A normal distribution with variance $\sigma^2$ is $\sigma$-sub-Gaussian; so is any bounded distribution. Sub-Gaussian means the chance of a very valuable plan falls off at least as fast as a normal tail: there are no rare, enormously valuable ideas. Whatever \Ber{} does, $U\le\bar U\le v$, because the realized action always lies in $\Rcal\subseteq\Acal(\xi)$.
\end{intuition}

\subsection{Three lemmas}

\begin{lemma}[Selection bound]\label{lem:selection}
Let $Z-\E[Z]$ be $\sigma$-sub-Gaussian, and let $A$ be any event with $P(A)=q\in(0,1]$. Then
\[
\E[Z\mid A]\ \le\ \E[Z]+\sigma\sqrt{2\ln(1/q)}.
\]
\end{lemma}

\begin{proof}
Let $X=Z-\E[Z]$ and fix $\theta>0$. Since $x\mapsto e^{\theta x}$ is convex, Jensen's inequality applied to the conditional distribution given $A$ gives
\[
e^{\theta\E[X\mid A]}\ \le\ \E\big[e^{\theta X}\mid A\big]=\frac{\E\big[e^{\theta X}\mathbf 1_A\big]}{q}\ \le\ \frac{\E\big[e^{\theta X}\big]}{q},
\]
where the last step uses $e^{\theta X}\ge0$. Taking logs and using the sub-Gaussian bound $\ln\E[e^{\theta X}]\le\theta^2\sigma^2/2$,
\[
\E[X\mid A]\ \le\ \frac{\theta\sigma^2}{2}+\frac{\ln(1/q)}{\theta}.
\]
If $q<1$, the right side is minimized at $\theta=\sqrt{2\ln(1/q)}/\sigma$, where it equals $\sigma\sqrt{2\ln(1/q)}$. If $q=1$, then $\E[X\mid A]=\E[X]=0$.
\end{proof}

\begin{intuition}
No matter how cleverly you select, and no matter what you know, the best a share $q$ of a thin-tailed population can average is about $\sigma\sqrt{2\ln(1/q)}$ above the mean. Becoming ten times more selective adds only a little to $\sqrt{\ln(1/q)}$. This is the mathematical core of every ceiling below: with thin tails, more selection has sharply diminishing returns.
\end{intuition}

\begin{lemma}[Screening bound]\label{lem:screening}
For any policy, $\Ycal\le\bar R\,\E[U^+]$, where $U^+=\max\{U,0\}$.
\end{lemma}

\begin{proof}
Each plan needs at least one signal, so $\E_\pi[N_\pi]\ge1$ and plans are examined at rate $R=n_A/(\tau_s\E_\pi[N_\pi])\le\bar R$. With $A=\{\accept\}$, $Q=R\,P(A)$, so
\[
\Ycal=Q\,\E[U\mid A]=R\,\E[U\mathbf 1_A]\ \le\ R\,\E[U^+]\ \le\ \bar R\,\E[U^+].\qedhere
\]
\end{proof}

\begin{lemma}[Output bound under selection]\label{lem:output}
Suppose $U\le Z$ for every plan, where $Z$ has mean $m_Z$ and $Z-m_Z$ is $\sigma$-sub-Gaussian, and suppose capacity satisfies $0<K\le\bar R/\sqrt e$. Then every policy with $Q(\pi)\le K$ has
\[
\Ycal\ \le\ B(K,\bar R)\equiv K\Big(m_Z^++\sigma\sqrt{2\ln(\bar R/K)}\Big).
\]
Moreover, $B$ is increasing in $\bar R$, and increasing in $K$ on $(0,\bar R/\sqrt e]$.
\end{lemma}

\begin{proof}
If $Q=0$, then $\Ycal=0\le B$. Otherwise let $R\le\bar R$ be the examination rate and $q=P(\accept)=Q/R\in(0,1]$. Since $U\le Z$, Lemma~\ref{lem:selection} applied to $Z$ gives
\[
\Ycal=Q\,\E[U\mid A]\le Q\,\E[Z\mid A]\le Q\big(m_Z+\sigma\sqrt{2\ln(1/q)}\big)\le Q\,m_Z^++\sigma f(Q,R),
\]
where $f(Q,R)=Q\sqrt{2\ln(R/Q)}$. It remains to bound $f$.
\begin{itemize}
  \item $f$ is increasing in $R$: $\partial f/\partial R=Q/\big(R\sqrt{2\ln(R/Q)}\big)>0$ when $Q<R$.
  \item Write $L=\ln(R/Q)$. Then $\partial f/\partial Q=\sqrt{2L}-1/\sqrt{2L}$, which is positive exactly when $L>1/2$, that is when $Q<R/\sqrt e$. So $f(\cdot,R)$ rises up to $Q=R/\sqrt e$, where it equals $R/\sqrt e$, and falls after.
\end{itemize}
Since $K\le\bar R/\sqrt e$, $f(\cdot,\bar R)$ is increasing on $(0,K]$. Two cases:
\begin{itemize}
  \item If $Q\le R/\sqrt e$: $f(Q,R)\le f(Q,\bar R)\le f(K,\bar R)$.
  \item If $Q>R/\sqrt e$: $f(Q,R)\le R/\sqrt e<Q\le K\le K\sqrt{2\ln(\bar R/K)}=f(K,\bar R)$, using $\ln(\bar R/K)\ge1/2$.
\end{itemize}
In both cases $f(Q,R)\le f(K,\bar R)$, and $Q\,m_Z^+\le K\,m_Z^+$, which gives $\Ycal\le B(K,\bar R)$.

For the monotonicity, $B$ is increasing in $\bar R$ by inspection. In $K$, with $L=\ln(\bar R/K)\ge1/2$,
\[
\frac{\partial B}{\partial K}=m_Z^++\sigma\Big(\sqrt{2L}-\frac{1}{\sqrt{2L}}\Big)\ \ge\ 0,
\]
because $\sqrt{2L}\ge1\ge1/\sqrt{2L}$.
\end{proof}

\begin{intuition}
Output is the number of plans pursued times their average value. The number is capped by capacity, and the average value of the top share of a thin-tailed pool is capped by Lemma~\ref{lem:selection}. The bound grows with capacity and with the size of the pool screened, but only logarithmically with the pool. Every ceiling below is this bound with a particular $Z$ and a particular cap on $K$.
\end{intuition}

%---------------------------------------------------------------------
\subsection{One improvement at a time}
%---------------------------------------------------------------------

\subsubsection*{Result 1: better problem solving alone has a ceiling}

\begin{proposition}\label{prop:ai}
Fix \Nap{} and \Lan. Along any path of better problem solving, for all $t$,
\[
\Ycal_t\ \le\ \frac{n_A}{\tau_s}\;\E_{\xi\sim F}\Big[\max\Big\{0,\ \max_{r\in\Rcal}W(r;\xi)\Big\}\Big]\ <\ \infty,
\]
so $g_{\Ycal}\le0$.
\end{proposition}

\begin{proof}
\begin{enumerate}
  \item Whatever \Ber{} does, the executed action lies in $\Rcal$, so $Y\le\bar U(\xi)$ and hence $U_t\le\bar U$ at every date.
  \item By Lemma~\ref{lem:screening}, $\Ycal_t\le\bar R\,\E[U_t^+]\le\bar R\,\E[\bar U^+]$. The right side is the displayed bound, and it does not depend on $t$, because $\bar R$ depends only on \Nap, and $\bar U$ only on the plans and on $\Rcal$.
  \item By~\eqref{eq:Ubar}, $\bar U\le v$, so $\E[\bar U^+]\le\E[v^+]\le\E|v|<\infty$.
  \item A bounded sequence satisfies $\frac1t\ln\Ycal_t\le\frac1t\ln C\to0$. \qedhere
\end{enumerate}
\end{proof}

\begin{intuition}
AI can do two things: let more plans through, up to the executor's capacity, and make each plan's details better, up to its best feasible value. Both have a ceiling that AI cannot move: one set by \Lan's capacity, the other by what $\Rcal$ allows. So AI raises output once and then stops.
\end{intuition}

\subsubsection*{Result 2: better ideation alone has a ceiling}

\begin{proposition}\label{prop:judgment}
Fix the signal time $\tau_s$ and \Lan, let $\bar U$ be $\sigma$-sub-Gaussian with mean $\bar m$, and let $K\le\mu_L\le\bar R/\sqrt e$. For any \Ber{} and any signal technology, including perfect information,
\[
\Ycal\ \le\ K\bigg(\max\Big\{0,\ \E_{\xi\sim F}\Big[\max_{r\in\Rcal}W(r;\xi)\Big]\Big\}+\sigma\sqrt{2\ln\frac{n_A}{K\tau_s}}\bigg).
\]
\end{proposition}

\begin{proof}
$U\le\bar U$ for every plan. Apply Lemma~\ref{lem:output} with $Z=\bar U$, $m_Z=\bar m=\E[\max_{r\in\Rcal}W(r;\xi)]$, and $\bar R=n_A/\tau_s$.
\end{proof}

\begin{intuition}
Better judgment lets \Nap{} pick the truly best of the plans he examines, but it cannot make the plans he examines better. Even with perfect information, the best he can do is sort them perfectly, and the top share of a thin-tailed pool is only so good. The ceiling depends on how many plans he can examine and how many he can pass on, and on nothing about his signals. Faster screening is the other half of better ideation; Result~5 shows that with thin tails it raises output only like $\sqrt t$.
\end{intuition}

\subsubsection*{Heavy tails change this}\label{sec:heavy}

\begin{proposition*}[Heavy tails]
Suppose expected values have a Pareto tail, $P(U>x)=(x/x_m)^{-\alpha}$ for $x\ge x_m>0$, with $\alpha>1$, and capacity $K$ is fixed. If screening speeds up at rate $g$, $\bar R_t=\bar R_0e^{gt}$, then $g_{\Ycal}\le g/\alpha$, with equality under perfectly informative signals.
\end{proposition*}

\begin{proof}
\begin{enumerate}
  \item \emph{Top-share mean.} The top share $q$ of $U$ lies above $c_q$ with $(c_q/x_m)^{-\alpha}=q$, i.e.\ $c_q=x_mq^{-1/\alpha}$. For a Pareto distribution, $\E[U\mid U\ge c]=\frac{\alpha}{\alpha-1}c$, so the mean of the top share is
  \[
  M(q)=\frac{\alpha}{\alpha-1}\,x_m\,q^{-1/\alpha}.
  \]
  \item \emph{No selection beats perfect sorting.} Let $A$ be any event with $P(A)=q$. Then $(U-c_q)(\mathbf 1\{U\ge c_q\}-\mathbf 1_A)\ge0$ pointwise, and $\E[\mathbf 1\{U\ge c_q\}-\mathbf 1_A]=0$, so $\E[U\mathbf 1_A]\le\E[U\mathbf 1\{U\ge c_q\}]$, that is $\E[U\mid A]\le M(q)$.
  \item \emph{Upper bound.} With $q=Q/R\ge Q/\bar R$ and $M$ decreasing, $\Ycal=Q\,\E[U\mid A]\le Q\,M(Q/\bar R)=\frac{\alpha}{\alpha-1}x_m\,Q^{1-1/\alpha}\bar R^{1/\alpha}$, which is increasing in $Q$, so $\Ycal_t\le K\,M(K/\bar R_t)\propto\bar R_t^{1/\alpha}=\bar R_0^{1/\alpha}e^{gt/\alpha}$.
  \item \emph{Attained.} With perfectly informative signals, one signal reveals $U$, so $\E[N]=1$, plans are examined at rate $\bar R_t$, and \Nap{} accepts exactly the top share $K/\bar R_t$. Since $U\ge x_m>0$, the cutoff is positive, so $\Ycal_t=K\,M(K/\bar R_t)$, which grows at exactly $g/\alpha$. \qedhere
\end{enumerate}
\end{proof}

\begin{intuition}
With a heavy tail, there are always rare ideas that are vastly better than the rest, and the more ideas you look at, the better the best one you find, by a power of the pool size rather than a logarithm. So faster screening can sustain growth, at rate $g/\alpha$. Lasting growth from better ideation needs both faster screening \emph{and} rare, enormously valuable plans. Section~\ref{sec:physical} shows that the physical frontier closes this way out in production.
\end{intuition}

\subsubsection*{Result 3: better execution alone has a ceiling}

\begin{proposition}\label{prop:execution}
Fix \Nap{} and \Ber, let $v(\xi)$ be $\sigma$-sub-Gaussian with mean $m$, and let $\mu_B\le\bar R/\sqrt e$. Along any path of better execution, for all $t$,
\[
\Ycal_t\ \le\ \mu_B\bigg(\max\big\{0,\ m\big\}+\sigma\sqrt{2\ln\frac{n_A}{\mu_B\tau_s}}\bigg).
\]
\end{proposition}

\begin{proof}
\begin{enumerate}
  \item Capacity is $K_t=\min\{\mu_B,\mu_{L,t}\}\le\mu_B\le\bar R/\sqrt e$, however fast \Lan{} becomes.
  \item Whatever \Lan's budget set, the executed action lies in $\Acal(\xi)$, so $U_t\le v$.
  \item Lemma~\ref{lem:output} with $Z=v$ gives $\Ycal_t\le B(K_t,\bar R)$, and $B$ is increasing in $K$ on $(0,\bar R/\sqrt e]$, so $\Ycal_t\le B(\mu_B,\bar R)$, which is the displayed bound. It does not depend on $t$. \qedhere
\end{enumerate}
\end{proof}

\begin{intuition}
Better hands cannot plan faster. However quickly \Lan{} can act and however far his menu reaches, only $\mu_B$ plans a year are worked out, so plans queue at \Ber. The bound uses the ideal value $v$ rather than the best feasible value, because the budget set itself changes along the path.
\end{intuition}

\paragraph{In the illustration.} Following the procedure of Section~\ref{sec:experiment}:
\begin{center}
\renewcommand{\arraystretch}{1.12}
\begin{tabular}{@{}l r r r r r l@{}}
\toprule
& Year 0 & 10 & 20 & 60 & Growth & Bottleneck\\
\midrule
No improvement & 1.57 & 1.57 & 1.57 & 1.57 & 0\% & \Ber\\
Better ideation alone & 1.57 & 4.02 & 6.38 & 12.75 & 0.9\% & \Ber\\
\textcolor{ber}{AI} & 1.57 & 4.14 & 4.43 & 4.47 & 0\% & \Lan\\
Better execution alone & 1.57 & 2.46 & 2.56 & 2.57 & 0\% & \Ber\\
\bottomrule
\end{tabular}
\end{center}
\begin{itemize}
  \item \textbf{AI} removes \Ber{} as the bottleneck (capacity rises from 1 to 2) and shrinks his losses towards zero. Output nearly triples, then stops at $4.47$, the output with $K=2$, $\kappa=1$, and year-0 screening: Result~1.
  \item \textbf{Better ideation} grows fast at first, as sharper signals pay off, then slows to under $1\%$ a year: sharper signals hit the full-information ceiling (Result~2), and faster screening helps only like $\sqrt t$ (Result~5).
  \item \textbf{Better execution} only removes the execution loss: \Ber{} was already the bottleneck, so capacity stays at one plan a year. Output stops at $2.57$: Result~3.
\end{itemize}

%---------------------------------------------------------------------
\subsection{Two improvements, not three}
%---------------------------------------------------------------------

Each pair of improvements leaves one agent behind, and that agent sets the ceiling:
\begin{center}
\renewcommand{\arraystretch}{1.12}
\begin{tabular}{@{}l l l l@{}}
\toprule
Pair & Left behind & What caps output & Results\\
\midrule
RSI & \Lan & the physical frontier and $\mu_L$ & 4--5\\
AI with robotics & \Nap & the speed of screening, $n_A/\tau_s$ & 6\\
Ideation with robotics & \Ber & the speed of solving, $\mu_B$ & 7\\
\bottomrule
\end{tabular}
\end{center}

\subsubsection{RSI: \texorpdfstring{\Lan}{Lannes} left behind}\label{sec:physical}

Results 2 and 3 leave one way out for RSI: heavy-tailed ideas. This subsection shows that once the physical frontier is taken into account, that way out is closed.

\paragraph{A one-dimensional budget set.} Let actions be real numbers, $\Acal=[\underline a,\overline a]$. \Ber{} can propose any $\hat a\in\Acal$, but \Lan's budget set is finite:
\[
\Rcal=\{r_1<r_2<\dots<r_J\},\qquad\tilde a=\arg\min_{r\in\Rcal}|r-\hat a|.
\]
$\Rcal$ is the set of physical technologies that exist. The \textbf{physical frontier} is $\bar r=\max\Rcal$: anything beyond it is carried out at $\bar r$.

\paragraph{Better plans sit further out.} Let a plan's ideal action measure its \textbf{ambition}, and let value rise with ambition:
\begin{equation}\label{eq:ambition}
v(\xi)=\beta\,a^*(\xi),\qquad W(a;\xi)=v(\xi)-\tfrac H2\big(a-a^*(\xi)\big)^2,\qquad\beta,H>0.
\end{equation}

\begin{intuition}
More valuable plans call for more extreme actions: a faster aircraft needs a stronger material, a bigger campaign needs more supply. Plotted against ambition $a^*$, potential value rises in a straight line. Realized value tracks it while a feasible technology is close by, dips between feasible points, peaks just beyond the frontier at $a^*=\bar r+\beta/H$, and then \emph{falls}: an idea that demands far more than any technology can deliver is worth less than a modest idea that can actually be done. Because \Nap's signals are about $U$, he learns about the realized curve, not the potential line, and does not chase plans beyond the frontier.
\end{intuition}

\paragraph{Result 4: the physical frontier bounds output.}

\begin{proposition}[The physical frontier bounds output]\label{prop:frontier}
Under~\eqref{eq:ambition}, for any \Nap, any \Ber, and any distribution of plans,
\[
U,\ Y\ \le\ \beta\bar r+\frac{\beta^2}{2H}\qquad\text{and}\qquad\Ycal\ \le\ \mu_L\Big(\beta\bar r+\frac{\beta^2}{2H}\Big).
\]
\end{proposition}

\begin{proof}
\begin{enumerate}
  \item The realized action is some $r\in\Rcal$, so $r\le\bar r$.
  \item For fixed $r$, the payoff as a function of the plan's ambition, $\beta a^*-\frac H2(r-a^*)^2$, is concave in $a^*$ with derivative $\beta+H(r-a^*)$, which is zero at $a^*=r+\beta/H$. Its maximum value is
  \[
  \beta\Big(r+\frac\beta H\Big)-\frac H2\Big(\frac\beta H\Big)^2=\beta r+\frac{\beta^2}{2H}.
  \]
  So $Y\le\beta r+\beta^2/(2H)\le\beta\bar r+\beta^2/(2H)$ for every plan, and the same holds for its expectation $U$.
  \item At most $\mu_L$ plans are carried out per unit of time, so $\Ycal=Q\,\E[Y\mid\accept]\le\mu_L\big(\beta\bar r+\beta^2/(2H)\big)$. \qedhere
\end{enumerate}
\end{proof}

\paragraph{Ambition and the frontier are complements.} The frontier does more than cap output: it interacts with the plan.

\begin{lemma}[Increasing differences]\label{lem:complements}
Under~\eqref{eq:ambition}, consider a plan whose solution lies at or beyond the frontier, $\hat a\ge\bar r$ (for example, $a^*>\bar r$ and a perfect \Ber). Its realized value is
\[
Y(a^*,\bar r)=\beta a^*-\tfrac H2\,(a^*-\bar r)^2,
\]
which has increasing differences in ambition and the frontier, and is maximized over ambition just beyond the frontier:
\[
\frac{\partial^2Y}{\partial a^*\,\partial\bar r}=H>0,\qquad\arg\max_{a^*}Y=\bar r+\frac{\beta}{H}.
\]
\end{lemma}

\begin{proof}
Since $\hat a\ge\bar r=\max\Rcal$, the point of $\Rcal$ nearest to $\hat a$ is $\bar r$, so $\tilde a=\bar r$ and $Y=W(\bar r;\xi)=\beta a^*-\frac H2(\bar r-a^*)^2$. Differentiating, $\partial Y/\partial\bar r=H(a^*-\bar r)$, so $\partial^2Y/\partial a^*\partial\bar r=H$. And $\partial Y/\partial a^*=\beta-H(a^*-\bar r)$, which is decreasing in $a^*$ and zero at $a^*=\bar r+\beta/H$.
\end{proof}

\begin{intuition}
Payoff is supermodular in the idea and the action that can be taken. The gain from moving the frontier out, $H(a^*-\bar r)$, is largest for the most ambitious plans, so better execution raises the return to better ideas; and the best plans to pursue sit a fixed distance $\beta/H$ beyond the frontier, so the frontier pulls ambition along with it. In the Gaussian example, where $W=v(\xi)+w(a)$, the cross-partial is zero: there better execution works only through $\kappa_L$ and capacity.
\end{intuition}

\paragraph{Result 5: with thin tails, RSI levels off.}

\begin{proposition}\label{prop:thin}
Fix \Lan, let $\bar U$ be $\sigma$-sub-Gaussian with mean $\bar m$, and let $\mu_L\le\bar R_0/\sqrt e$. Suppose signals speed up at most exponentially, $\tau_{s,t}\ge\tau_{s,0}e^{-gt}$, and \Nap's technology only improves over time. Then for any paths of better problem solving and better ideation,
\[
\Ycal_t\ \le\ \frac{n_L}{\tau_L}\bigg(\max\Big\{0,\ \E_{\xi\sim F}\Big[\max_{r\in\Rcal}W(r;\xi)\Big]\Big\}+\sigma\sqrt{2\ln\frac{n_A\tau_L}{n_L\tau_{s,0}}+2gt}\bigg).
\]
In particular output grows at most like $\sqrt t$, and $g_{\Ycal}\le0$.
\end{proposition}

\begin{proof}
\begin{enumerate}
  \item At each date $t$, $U_t\le\bar U$, and $\bar U$ does not change over time, because it depends only on $F$ and $\Rcal$.
  \item Capacity is $K_t=\min\{\mu_{B,t},\mu_L\}\le\mu_L\le\bar R_0/\sqrt e\le\bar R_t/\sqrt e$, since $\tau_{s,t}\le\tau_{s,0}$.
  \item Lemma~\ref{lem:output} with $Z=\bar U$ gives $\Ycal_t\le B(K_t,\bar R_t)\le B(\mu_L,\bar R_t)$, by monotonicity in $K$.
  \item By monotonicity in $\bar R$ and $\bar R_t\le\bar R_0e^{gt}$, $\ln(\bar R_t/\mu_L)\le\ln(\bar R_0/\mu_L)+gt$, which gives the displayed bound.
  \item A bound of order $\sqrt t$ satisfies $\frac1t\ln\Ycal_t\le\frac1t\ln(c_1+c_2\sqrt t)\to0$. \qedhere
\end{enumerate}
\end{proof}

\begin{intuition}
Exponentially faster screening means \Nap{} examines exponentially more plans, but with thin tails the best of $e^{gt}$ plans improves only like $\sqrt{2gt}$. Better solving cannot help, because it cannot lift plans above their best feasible value; better judgment cannot help, by Result~2. So AI, better ideation, and RSI all level off. Their levels differ, possibly by a lot, but their long-run growth rate is zero.
\end{intuition}

\begin{corollary}[Without better execution, nothing grows]\label{cor:frontier}
Under~\eqref{eq:ambition}, if $\bar r_t\le\bar r_0e^{g_Rt}$ and $\mu_{L,t}\le\mu_{L,0}e^{g_Lt}$ with $\bar r_0>0$, then for any paths of better problem solving and better ideation,
\[
\limsup_{t\to\infty}\tfrac1t\ln\Ycal_t\ \le\ g_R+g_L.
\]
Without better execution ($g_R=g_L=0$), $g_{\Ycal}\le0$ under AI, better ideation, or RSI, even with heavy-tailed plans.
\end{corollary}

\begin{proof}
By Proposition~\ref{prop:frontier}, $\Ycal_t\le\mu_{L,0}e^{g_Lt}\big(\beta\bar r_0e^{g_Rt}+\beta^2/(2H)\big)\le\mu_{L,0}\big(\beta\bar r_0+\beta^2/(2H)\big)e^{(g_L+g_R)t}$, using $g_R\ge0$ (if $g_R<0$ the bound $\mu_{L,0}e^{g_Lt}(\beta\bar r_0+\beta^2/(2H))$ applies instead). Taking logs, dividing by $t$, and letting $t\to\infty$ gives the result.
\end{proof}

\begin{intuition}
Heavy-tailed potential value does not help if the tail lies beyond what can be done. Whatever the strategist dreams up and however precisely the solver works it out, what reaches the world is capped by the best available technology.
\end{intuition}

\subsubsection{AI with robotics: \texorpdfstring{\Nap}{Napoleon} left behind}

\paragraph{Result 6: AI with robotics reaches \Nap's ceiling.}

\begin{proposition}\label{prop:airobotics}
Fix \Nap. Along any paths of better problem solving and better execution, for all $t$,
\[
\Ycal_t\ \le\ \frac{n_A}{\tau_s}\;\E_{\xi\sim F}\big[\max\{0,\ v(\xi)\}\big],
\]
and once $\min\{\mu_{B,t},\mu_{L,t}\}\ge n_A/\tau_s$, \Nap's capacity constraint is slack for every policy.
\end{proposition}

\begin{proof}
\begin{enumerate}
  \item The executed action lies in $\Rcal\subseteq\Acal(\xi)$, so $U_t\le v$, whatever \Ber{} and \Lan{} do.
  \item By Lemma~\ref{lem:screening}, $\Ycal_t\le\bar R\,\E[U_t^+]\le\bar R\,\E[v^+]$, with $\bar R=n_A/\tau_s$ fixed. The bound is finite by Assumption~\ref{ass:results}(b).
  \item Every policy has $Q(\pi)\le R(\pi)\le n_A/\tau_s$, since each plan needs at least one signal. If $K_t=\min\{\mu_{B,t},\mu_{L,t}\}\ge n_A/\tau_s$, then $Q(\pi)\le K_t$ for every policy, so the constraint in Problem N is slack. \qedhere
\end{enumerate}
\end{proof}

\begin{intuition}
With both downstream stages improving, capacity eventually outruns \Nap's screening. He stops rationing and pursues every plan he expects to be worth something; his own screening time becomes the binding constraint. Output rises to a much higher level than under AI alone, but the ceiling depends only on \Nap.
\end{intuition}

\subsubsection{Ideation with robotics: \texorpdfstring{\Ber}{Berthier} left behind}

\paragraph{Result 7: ideation with robotics is capped by \Ber.}

\begin{proposition}\label{prop:ideationrobotics}
Fix \Ber, let $v(\xi)$ be $\sigma$-sub-Gaussian with mean $m$, and let $\mu_B\le\bar R_0/\sqrt e$. Suppose signals speed up at most exponentially, $\tau_{s,t}\ge\tau_{s,0}e^{-gt}$, and \Nap's technology only improves over time. Then for any paths of better ideation and better execution,
\[
\Ycal_t\ \le\ \mu_B\bigg(\max\{0,\ m\}+\sigma\sqrt{2\ln\frac{n_A}{\mu_B\tau_{s,0}}+2gt}\bigg),
\]
so output grows at most like $\sqrt t$, and $g_{\Ycal}\le0$.
\end{proposition}

\begin{proof}
\begin{enumerate}
  \item Capacity is $K_t=\min\{\mu_B,\mu_{L,t}\}\le\mu_B\le\bar R_0/\sqrt e\le\bar R_t/\sqrt e$, since $\tau_{s,t}\le\tau_{s,0}$.
  \item $U_t\le v$ at every date, whatever \Lan's budget set.
  \item Lemma~\ref{lem:output} with $Z=v$ gives $\Ycal_t\le B(K_t,\bar R_t)\le B(\mu_B,\bar R_t)\le B(\mu_B,\bar R_0e^{gt})$, by monotonicity in $K$ and in $\bar R$. This is the displayed bound.
  \item A bound of order $\sqrt t$ implies $g_{\Ycal}\le0$. \qedhere
\end{enumerate}
\end{proof}

\begin{intuition}
Without better problem solving, good plans pile up unsolved. Better ideation finds better plans and better execution can carry out whatever is solved, but only $\mu_B$ plans a year get worked out, so growth would have to come from ever-better ideas, and with thin tails those run out.
\end{intuition}

\paragraph{In the illustration.}
\begin{center}
\renewcommand{\arraystretch}{1.12}
\begin{tabular}{@{}l r r r r r l@{}}
\toprule
& Year 0 & 10 & 20 & 60 & Growth & Bottleneck\\
\midrule
\textcolor{nap}{RSI} & 1.57 & 8.95 & 13.98 & 27.12 & 0.9\% & \Lan\\
\textcolor{lan}{AI with robotics} & 1.57 & 11.75 & 13.99 & 14.27 & 0\% & \Nap\\
Ideation with robotics & 1.57 & 4.91 & 7.37 & 13.75 & 0.8\% & \Ber\\
\bottomrule
\end{tabular}
\end{center}
\begin{itemize}
  \item Each pair is capped by the agent it leaves behind: \Lan{} for RSI, \Nap{} for AI with robotics, \Ber{} for ideation with robotics.
  \item Improvements are complements. AI alone adds $2.90$ to year-60 output and better execution alone adds $1.00$, but together they add $12.70$: each removes the other's bottleneck.
  \item AI with robotics in production ends at exactly the level of AI in research ($14.27$, Section~\ref{sec:research}): once execution is no longer scarce, production behaves like research.
\end{itemize}

%---------------------------------------------------------------------
\subsection{All three improvements}
%---------------------------------------------------------------------

\paragraph{Result 8: all three improvements sustain growth.}

\begin{proposition}\label{prop:allthree}
Take Gaussian plans, $v(\xi)\sim N(m,\sigma^2)$ with the sieve closed under adding constants, and Gaussian signals. Suppose better ideation speeds up screening at rate $g\ge0$, $\tau_{s,t}=\tau_{s,0}e^{-gt}$, without making signals noisier; better problem solving raises capacity at rate $g_B\ge0$, $\mu_{B,t}=\mu_{B,0}e^{g_Bt}$; and better execution raises capacity at rate $g_L\ge0$, $\mu_{L,t}=\mu_{L,0}e^{g_Lt}$; and neither loss rises over time. Then
\[
g_{\Ycal}=\lim_{t\to\infty}\tfrac1t\ln\Ycal_t=\min\{g,\ g_B,\ g_L\}.
\]
\end{proposition}

\begin{proof}
Write $\bar R_t=n_A/\tau_{s,t}=\bar R_0e^{gt}$ and $K_t=\min\{\mu_{B,t},\mu_{L,t}\}$, which grows at rate $g_K=\min\{g_B,g_L\}$ in the sense that $\frac1t\ln K_t\to g_K$. Whatever the budget set, $U_t\le v$, and $v$ is $\sigma$-sub-Gaussian around $m$.

\emph{Upper bound.} There are two cases.
\begin{itemize}
  \item If $g\le g_K$: by Lemma~\ref{lem:screening}, $\Ycal_t\le\bar R_t\,\E[U_t^+]\le\bar R_t\,\E[v^+]$, which grows at rate $g=\min\{g,g_K\}$.
  \item If $g_K<g$: the ratio $\bar R_t/K_t$ tends to infinity, so eventually $K_t\le\bar R_t/\sqrt e$. Lemma~\ref{lem:output} with $Z=v$ gives
  \[
  \Ycal_t\ \le\ K_t\Big(m^++\sigma\sqrt{2\ln(\bar R_t/K_t)}\Big)=K_t\,O(\sqrt t),
  \]
  which grows at rate $g_K=\min\{g,g_K\}$.
\end{itemize}
In both cases $\limsup_t\frac1t\ln\Ycal_t\le\min\{g,g_K\}$.

\emph{Lower bound.} Output is the maximum of Problem N, so it is at least the output of any feasible policy. Consider the policy that takes one signal per plan and accepts if $\mu_1\ge\max\{c,z_t\}$, for a fixed $c>0$ and a capacity cutoff $z_t$ chosen so that the accepted flow does not exceed $K_t$.
\begin{itemize}
  \item \emph{Value of accepted plans.} Since $\E[U\mid s_1]=\mu_1$ and accepted plans have $\mu_1\ge c$, iterated expectations give $\E[U\mid\accept]\ge c$.
  \item \emph{Acceptance probability.} By the Gaussian-plans lemma, $\mu_1\sim N(m_{U,t},\rho_{1,t}\sigma^2)$ with $m_{U,t}=m-\kappa_{B,t}-\kappa_{L,t}$ and $\rho_{1,t}=1/(1+r_t)$. If $c\le m_{U,t}$, then $P(\mu_1\ge c)\ge1/2$. Otherwise
  \[
  P(\mu_1\ge c)=1-\Phi\Big(\frac{c-m_{U,t}}{\sigma\sqrt{\rho_{1,t}}}\Big),
  \]
  which is increasing in $m_{U,t}$ and in $\rho_{1,t}$. Both weakly rise over time, since neither loss rises and signals do not get noisier. So $P(\mu_1\ge c)\ge p\equiv\min\{p_0,1/2\}>0$ for all $t$, where $p_0$ is the value at $t=0$.
  \item \emph{Flow.} One signal per plan means plans are examined at rate $\bar R_t$, so the accepted flow is
  \[
  \min\{\bar R_t\,P(\mu_1\ge c),\ K_t\}\ \ge\ \min\{\bar R_t\,p,\ K_t\}.
  \]
\end{itemize}
So $\Ycal_t\ge c\,\min\{\bar R_0\,p\,e^{gt},\ K_t\}$, and $\liminf_t\frac1t\ln\Ycal_t\ge\min\{g,g_K\}$. Combining the two bounds, the limit exists and equals $\min\{g,g_B,g_L\}$.
\end{proof}

\begin{intuition}
Output is the number of plans pursued times their average value. With all three improving, every constraint on the number keeps relaxing: \Nap{} screens more plans, \Ber{} solves more, and \Lan{} carries out more. The number of projects grows exponentially, at the pace of whichever improvement is slowest, and that growth does not depend on finding ever-better ideas, so it survives thin tails. It is a weakest link. Leave any improvement out and its rate is zero, so $g_{\Ycal}=0$: Results 1--7 again, in the Gaussian setting.
\end{intuition}

\paragraph{In the illustration: the full picture.}
\begin{center}
\renewcommand{\arraystretch}{1.12}
\begin{tabular}{@{}l c c c r r l@{}}
\toprule
& \textcolor{nap}{Ide.} & \textcolor{ber}{Sol.} & \textcolor{lan}{Exe.} & Year 60 & Growth & Bottleneck\\
\midrule
No improvement & & & & 1.57 & 0\% & \Ber\\
Better ideation alone & $\checkmark$ & & & 12.75 & 0.9\% & \Ber\\
\textcolor{ber}{AI} & & $\checkmark$ & & 4.47 & 0\% & \Lan\\
Better execution alone & & & $\checkmark$ & 2.57 & 0\% & \Ber\\
\midrule
\textcolor{nap}{RSI} & $\checkmark$ & $\checkmark$ & & 27.12 & 0.9\% & \Lan\\
\textcolor{lan}{AI with robotics} & & $\checkmark$ & $\checkmark$ & 14.27 & 0\% & \Nap\\
Ideation with robotics & $\checkmark$ & & $\checkmark$ & 13.75 & 0.8\% & \Ber\\
\midrule
\textbf{RSI with robotics} & $\checkmark$ & $\checkmark$ & $\checkmark$ & \textbf{9{,}654} & \textbf{10.1\%} & \Nap, improving\\
\bottomrule
\end{tabular}
\end{center}
With all three, output grows at about $10\%$ a year, the screening rate, which is the slowest of the three improvements ($1.1$ against $1.25$ for the other two), exactly as Result~8 predicts.

%=====================================================================
\section{A research sector}\label{sec:research}
%=====================================================================

So far every plan had to be carried out physically. Much valuable work does not: proving a theorem, designing a circuit, writing software, running a simulation. In such work, once the details are worked out, the job is done. We call this the \textbf{research sector}: \Nap{} picks the question and \Ber's solution \emph{is} the output.

\subsection{Setup}

The research sector is the same model with no executor. Formally, take $\Rcal=\Acal(\xi)$ (every action is feasible) and $\mu_L=\infty$ (no execution capacity limit). Then $\tilde a=\hat a$, $\kappa_L=0$, and
\[
Y=W(\hat a;\xi),\qquad U=v(\xi)-\kappa_B(\xi),\qquad \bar U(\xi)=v(\xi),\qquad K=\mu_B=\frac{n_B}{\tau_B}.
\]
The best feasible value is the ideal value itself, because nothing restricts which actions can be taken. Downstream capacity is \Ber's capacity alone, and better execution has nothing to act on.

\subsection{AI relaxes Napoleon's constraint}

\Nap's capacity constraint is $Q(\pi)\le K$. In production, $K=\min\{\mu_B,\mu_L\}\le\mu_L$: AI can lift capacity only until \Lan{} binds. In research, $K=\mu_B$, so AI lifts capacity without limit. On the other side, every policy satisfies
\[
Q(\pi)\ \le\ R(\pi)=\frac{n_A}{\tau_s\,\E_\pi[N_\pi]}\ \le\ \frac{n_A}{\tau_s},
\]
because each plan needs at least one signal. So once $\mu_B\ge n_A/\tau_s$, the capacity constraint can never bind.

\paragraph{Result 9: AI lifts research to \Nap's ceiling.}

\begin{proposition}\label{prop:resceiling}
In research, fix \Nap. Along any path of better problem solving, for all $t$,
\[
\Ycal_t\ \le\ \frac{n_A}{\tau_s}\;\E_{\xi\sim F}\big[\max\{0,\ v(\xi)\}\big],
\]
and once $\mu_{B,t}\ge n_A/\tau_s$, \Nap's capacity constraint is slack for every policy.
\end{proposition}

\begin{proof}
This is Proposition~\ref{prop:airobotics} with $\mu_{L,t}=\infty$: in research $U=v-\kappa_B\le v$ because $\kappa_B\ge0$ by~\eqref{eq:kappaB}; Lemma~\ref{lem:screening} gives the bound; and $K_t=\mu_{B,t}\ge n_A/\tau_s\ge Q(\pi)$ makes the constraint slack.
\end{proof}

\begin{intuition}
When solving becomes cheap and fast, \Nap{} no longer has to ration. He stops being picky: he accepts every plan he expects to be worth something, and he spends fewer signals on each plan, because the point of investigating was to choose among plans competing for scarce slots. His binding constraint becomes his own screening time. In production, AI with robotics reaches the same ceiling; AI alone cannot, because \Lan's capacity binds first.
\end{intuition}

\paragraph{Result 10: in research, RSI is enough.}

\begin{proposition}\label{prop:resgrowth}
In research with Gaussian plans, suppose better problem solving raises capacity at rate $g_B$, $\mu_{B,t}=\mu_{B,0}e^{g_Bt}$, and does not raise $\kappa_B$; and better ideation speeds up screening at rate $g$, $\tau_{s,t}=\tau_{s,0}e^{-gt}$, without making signals noisier. Then
\[
g_{\Ycal}=\lim_{t\to\infty}\tfrac1t\ln\Ycal_t=\min\{g_B,\ g\}.
\]
\end{proposition}

\begin{proof}
Research is the case $\Rcal=\Acal(\xi)$, $\kappa_L=0$, and $\mu_L=\infty$, so $K_t=\mu_{B,t}$ and $m_{U,t}=m-\kappa_{B,t}$. The proof of Proposition~\ref{prop:allthree} applies word for word with $K_t=\mu_{B,t}$ and $g_K=g_B$: the upper bound uses only $U_t\le v$, Lemma~\ref{lem:screening}, and Lemma~\ref{lem:output}; the lower bound uses only that $m_{U,t}$ and $\rho_{1,t}$ do not fall. So $g_{\Ycal}=\min\{g,g_B\}$.
\end{proof}

\begin{intuition}
With no \Lan{} to improve, RSI is already ``all three''. Research output is the number of projects completed times their average value; with RSI both constraints on the number keep relaxing, so research grows at the rate of the slower one, even with thin tails. In production the same improvements leave \Lan{} behind, so growth would have to come from ever-better ideas, and with thin tails those run out.
\end{intuition}

\subsection{Research versus production in the illustration}

Take the experiment of Section~\ref{sec:experiment}, but drop \Lan: $\kappa_L=0$ and no execution capacity. Year-0 research output is $2.57$ (versus $1.57$ in production, the difference being the absence of the execution loss).
\begin{center}
\renewcommand{\arraystretch}{1.12}
\begin{tabular}{@{}l r r r r@{}}
\toprule
& \multicolumn{2}{c}{Production} & \multicolumn{2}{c}{Research}\\
\cmidrule(lr){2-3}\cmidrule(lr){4-5}
& Year 60 & Growth & Year 60 & Growth\\
\midrule
\textcolor{ber}{AI} & 4.47 & 0\% & 14.27 & 0\%\\
\textcolor{lan}{AI with robotics} & 14.27 & 0\% & 14.27 & 0\%\\
Better ideation alone & 12.75 & 0.9\% & 13.75 & 0.8\%\\
\textcolor{nap}{RSI} & 27.12 & 0.9\% & 9{,}654 & 10.1\%\\
\bottomrule
\end{tabular}
\end{center}
\begin{itemize}
  \item \textbf{AI}: research output rises $5.6$-fold (from $2.57$ to $14.27$), against $2.8$-fold in production. \Nap{} goes from accepting $1$ plan a year to $10$, from $3$ signals per plan to $1$, and his cutoff falls to zero. The ceiling of $14.27$ is set by his screening speed alone.
  \item \textbf{Robotics} has nothing to do in research, so AI with robotics in research is just AI; in production it ends at the same $14.27$.
  \item \textbf{RSI}: research output grows at about $10\%$ a year, the screening rate. In production the same improvements give under $1\%$; adding better execution (RSI with robotics) restores the $10\%$.
\end{itemize}

%=====================================================================
\section{Discussion}
%=====================================================================

\subsection{What the model says}

The three kinds of work have three kinds of ceiling. Solving has a ceiling at the plan's best feasible value. Judgment has a ceiling at perfect information. Screening more ideas has no hard ceiling, but with thin tails its returns are logarithmic. And all of them are bounded by what can physically be done.

Organized by how many agents improve, the results say one thing: \emph{the agent left behind sets the limit}. One improvement raises output to a ceiling set by another agent. Two improvements raise it further, to a ceiling set by the third. Only improving all three sustains growth, at the pace of the slowest improvement. Research is the exception that proves the rule: with no physical stage there are only two agents, so RSI is enough. A natural reading is that AI's largest effects should appear first in research, and reach production only to the extent that execution improves too.

\subsection{What the model assumes, and how it could change}

\begin{itemize}
  \item \textbf{Ideas are free and the pool is fixed.} All of \Nap's time goes into evaluating ideas, none into producing them, and the pool $F$ never improves. If better ideation could shift $F$, not only improve signals, it could sustain growth through a different channel.
  \item \textbf{No recall, one idea at a time.} A rejected plan never returns. A more realistic strategist would hold many ideas at once, spend more signals on promising ones, and pick the best of the lot: a ranking-and-selection problem. The bounds in Section~\ref{sec:results} still apply, because Lemma~\ref{lem:selection} holds for any selection rule.
  \item \textbf{The solver uses the best approximation in his sieve.} In practice he would fit $\hat W$ from finitely many evaluations of $W$, adding estimation error and tying his solving time $\tau_B$ to the number of evaluations.
  \item \textbf{Losses do not depend on plan complexity.} If ambitious plans are harder to solve, a better solver makes high-value plans feasible and partly acts like better ideation.
  \item \textbf{The physical frontier is exogenous.} If AI helps design better technologies, it moves $\bar r$ indirectly. Whether real progress in AI is mostly better problem solving, better ideation, or a contributor to better execution is the empirical question on which the growth conclusions turn.
\end{itemize}

\subsection{Next steps}

\begin{enumerate}
  \item \textbf{Endogenous complexity.} Do ambitious plans also need more \Lan{} time, $\tau_L(\xi)$, or more \Ber{} time, $\tau_B(\xi)$? Lemma~\ref{lem:complements} is a first step: it makes the value of execution depend on the plan.
  \item \textbf{Measurement.} Estimate the tail of plan values and the pace of physical progress, and separate failures of judgment, solving, and execution in data.
  \item \textbf{Occupational choice and wages.} How does cheaper \Ber{} labour change demand for \Nap- and \Lan-like workers?
  \item \textbf{General equilibrium.} Let the numbers of each agent, $n_A$, $n_B$, $n_L$, respond to prices.
\end{enumerate}

\end{document}
